\BourbakiExerciseGroup

\begin{enumerate}
\def\labelenumi{\arabic{enumi})}
\tightlist
\item
  \begin{enumerate}
  \def\labelenumii{\alph{enumii})}
  \tightlist
  \item
    Let \(q\) be a power of a prime number. Show that for every integer \(n > 1\) , if \(h_i\) ( \(1 \leqslant i \leqslant r\) ) are the distinct prime divisors of the integer \(n\) , then the number of elements \(\zeta \in \mathbf{F}_{q^n}\) such that \(\mathbf{F}_{q^n} = \mathbf{F}_q(\zeta)\) is equal to
  \end{enumerate}
\end{enumerate}

\[
\nu = q ^ {n} - \sum_ {i} q ^ {n / h _ {i}} + \sum_ {i <   j} q ^ {n / h _ {i} h _ {j}} - \sum_ {i <   j <   k} q ^ {n / h _ {i} h _ {j} h _ {k}} + \dots + (- 1) ^ {r} q ^ {n / h _ {1} h _ {2} \dots h _ {r}}
\]

(observe that such an element is characterized by the property of not belonging to any of the fields \(F_{q}^{n/h_{i}}\) ). If \(h_{1} \leqslant \cdots \leqslant h_{r}\) , we have

\[
q ^ {n} - \sum_ {i} q ^ {n / h _ {i}} \leqslant \nu \leqslant q ^ {n} - q ^ {n / h _ {1}}.
\]

Examine the case where \(n\) is a prime power.

\begin{enumerate}
\def\labelenumi{\alph{enumi})}
\setcounter{enumi}{1}
\item
  Let \(b_{n}(q)\) be the number of monic irreducible polynomials of degree \(n\) over \(\mathbf{F}_q\) . Deduce from \(a\) ) an expression for \(b_{n}(q)\) . In particular, if \(q \geqslant 3\) , show that the number of polynomials (not necessarily monic) of degree \(n\) in \(\mathbf{F}_q[\mathbf{X}]\) which are irreducible is at least equal to \(q^n (q - 2) / n \geqslant q^{n + 1} / 3n\) .
\item
  Prove that \(\sum_{d\mid n}db_d(q) = q^n\) (establish the identity \(1 / (1 - q\mathrm{T}) = \prod_{n}(1 - \mathrm{T}^n)^{-b_n(q)}\) ).
\end{enumerate}

* Using the Möbius inversion formula (Lie, II, p.~176, Appendix), show that

\[
n b _ {n} (q) = \sum_ {d \mid n} \mu (n / d) q ^ {d}.
\]

Compare this result with that obtained in \(b\) . *

\begin{enumerate}
\def\labelenumi{\alph{enumi})}
\setcounter{enumi}{3}
\tightlist
\item
  Let \(a_{n}(q)\) be the number of monic polynomials without multiple factors in \(\mathbf{F}_q[\mathbf{X}]\) . Show that \(a_0(q) = 1, \quad a_1(q) = q, \quad a_n(q) = q^n - q^{n-1}\) for \(n > 1\) (establish the identity \(q^n = \sum_i a_{n-2i}(q) q^i\) ).
\end{enumerate}

\begin{enumerate}
\def\labelenumi{\arabic{enumi})}
\setcounter{enumi}{1}
\item
  Show that the number of monic polynomials of degree \(n\) of \(\mathbf{F}_q[\mathbf{X}]\) which have no root in \(\mathbf{F}_q\) is \(q^{n - q}(q - 1)^q\) for \(n \geqslant q\) (observe that every function \(\mathbf{F}_q \to \mathbf{F}_q\) is a polynomial function).
\item
  Let \(\mathbf{K}\) be a finite field of \(q\) elements; for every prime number \(l\) let \(N_{l}\) be the union of the extensions of \(\mathbf{K}\) (contained in an algebraic closure \(\Omega\) of \(\mathbf{K}\) ) whose degree is a power of \(l\) . Show that \(\mathbf{N}_{l}\) is an abelian extension of \(\mathbf{K}\) whose Galois group is isomorphic to the group \(\mathbf{Z}_{l}\) of \(l\) -adic integers. Deduce that \(\Omega\) is isomorphic to the tensor product \(\otimes \mathbf{N}_{l}\) , and hence,
\end{enumerate}

that \(\operatorname{Gal}(\Omega / \mathbf{K})\) is isomorphic to the product \(\prod_{l} \mathbf{Z}_l\) .

* 4) Decompose the polynomial \(\mathbf{X}^4 + 1 \in \mathbf{F}_p[\mathbf{X}]\) into irreducible factors. *

\begin{enumerate}
\def\labelenumi{\arabic{enumi})}
\setcounter{enumi}{4}
\tightlist
\item
  Let \(\mathbf{K}\) be a finite field of \(q\) elements where \(q\) is not a power of 2, and let \(a_1, a_2, b\) be three elements of \(\mathbf{K}\) such that \(a_1a_2 \neq 0\) . Show that the number \(\nu\) of solutions \((x_1, x_2) \in \mathbf{K}^2\) of the equation \(a_1x_1^2 + a_2x_2^2 = b\) is given by the following formulae:
\end{enumerate}

\(1^{\circ}\) if \(b = 0\) and \(-a_{1}a_{2}\) is not a square in \(\mathbf{K},\nu = 1\)

\(2^{\circ}\) if \(b \neq 0\) and \(-a_{1}a_{2}\) is not a square in \(\mathbf{K}\) , \(\nu = q + 1\) ;

\(3^{\circ}\) if \(b = 0\) and \(-a_{1}a_{2}\) is a square in \(\mathbf{K}\) , \(\nu = 2q - 1\) ;

\(4^{\circ}\) if \(b \neq 0\) and \(-a_{1}a_{2}\) is a square in K, \(\nu = q - 1\) .

(When \(-a_{1}a_{2}\) is not a square, adjoin to K a root of \(\mathbf{X}^{2} + a_{1}a_{2}\) and use the fact that the multiplicative group of a finite field is cyclic.)

\begin{enumerate}
\def\labelenumi{\arabic{enumi})}
\setcounter{enumi}{5}
\tightlist
\item
  Let \(\mathbf{K}\) be a finite field of \(q\) elements where \(q\) is not a power of 2.
\end{enumerate}

\begin{enumerate}
\def\labelenumi{\alph{enumi})}
\tightlist
\item
  Let \(a_{1}, a_{2}, \ldots, a_{2m}, b\) be elements of K such that \(a_{1}a_{2} \ldots a_{2m} \neq 0\) . Show that the number \(\nu\) of solutions \((x_{1}, \ldots, x_{2m}) \in \mathbf{K}^{2m}\) of the equation \(a_{1}x_{1}^{2} + \cdots + a_{2m}x_{2m}^{2} = b\) is given by the following formulae:
\end{enumerate}

\(1^{\circ}\) if \(b = 0\) and \((-1)^{m}a_{1}a_{2}\ldots a_{2m}\) is not a square in K, \(\nu = q^{2m - 1} - q^{m} + q^{m - 1}\) ;

\(2^{\circ}\) if \(b \neq 0\) and \((-1)^{m}a_{1}a_{2} \ldots a_{2m}\) is not a square in K, \(\nu = q^{2m-1} + q^{m-1}\) ;

\(3^{\circ}\) if \(b = 0\) and \((-1)^{m}a_{1}a_{2}\ldots a_{2m}\) is a square in K, \(\nu = q^{2m - 1} + q^{m} - q^{m - 1}\) ;

\(4^{\circ}\) if \(b \neq 0\) and \((-1)^{m}a_{1}a_{2} \ldots a_{2m}\) is a square in K, \(\nu = q^{2m - 1} - q^{m - 1}\) .

(Argue by induction on \(m\) , using Ex. 5.)

\begin{enumerate}
\def\labelenumi{\alph{enumi})}
\setcounter{enumi}{1}
\tightlist
\item
  Let \(a_1, a_2, \ldots, a_{2m+1}\) , \(b\) be elements of \(K\) such that \(a_1a_2 \ldots a_{2m+1} \neq 0\) . Show that the number \(\nu\) of solutions \((x_1, x_2, \ldots, x_{2m+1}) \in K^{2m+1}\) of the equation
\end{enumerate}

\[
a _ {1} x _ {1} ^ {2} + \dots + a _ {2 m + 1} x _ {2 m + 1} ^ {2} = b
\]

is given by the following formulae :

\(1^{\circ}\) if \(b = 0, \nu = q^{2m}\) ;

\(2^{\circ}\) if \(b \neq 0\) and \((-1)^{m}a_{1}a_{2} \ldots a_{2m+1}b\) is not a square in K, \(\nu = q^{2m} - q^{m}\) ;

\(3^{\circ}\) if \(b \neq 0\) and \((-1)^{m}a_{1}a_{2} \ldots a_{2m+1}b\) is a square in \(\mathbf{K}\) , \(\nu = q^{2m} + q^{m}\) . (Reduce to the case \(a\) ).)

\begin{enumerate}
\def\labelenumi{\arabic{enumi})}
\setcounter{enumi}{6}
\tightlist
\item
  Let K be a finite field, E an algebraic extension of degree n of K and \((\alpha_{1}, \ldots, \alpha_{n})\) a normal basis of E over K; suppose that the cyclic Galois group \(\operatorname{Gal}(\operatorname{E}/\operatorname{K})\) is generated by the cyclic permutation \(\sigma = (\alpha_{1}\alpha_{2} \ldots \alpha_{n})\) . Let \(\tau\) be an automorphism of the vector K-space E such that for each \(x \in E\) , \(\tau(x)\) is conjugate to x over E (cf.~V, p.~154, Ex. 8).
\end{enumerate}

\begin{enumerate}
\def\labelenumi{\alph{enumi})}
\item
  Show that if \(\mathbf{K}\) has at least 3 elements, \(\tau\) is a K-automorphism of the field E (reduce to the case where \(\tau(\alpha_1) = \alpha_1\) and consider \(\tau(\alpha_1 + c\alpha_j)\) for \(c \in \mathbf{K}\) ).
\item
  If \(K = F_{2}\) , show that for each element \(\rho \in \operatorname{Gal}(E/K)\) , the relations \(\rho(\alpha_{r}) = \alpha_{r'}\) and \(\rho(\alpha_{s}) = \alpha_{s'}\) imply \(r' - r \equiv s' - s \pmod{n}\) . Deduce that if n \textgreater{} 5, \(\tau\) is a K-automorphism of the field E. (Note that if \(\tau(\alpha_{1}) = \alpha_{1}\) and \(\tau(\alpha_{1} + \alpha_{h}) = \alpha_{1} + \alpha_{k}\) with \(k \neq h\) , we must have \(h + k = n + 2\) ; deduce from this that necessarily \(\tau(\alpha_{2}) = \alpha_{n}\) and \(\tau(\alpha_{4}) = \alpha_{4}\) or \(\tau(\alpha_{4}) = \alpha_{n-2}\) which leads to a contradiction.)
\item
  For \(\mathbf{K} = \mathbf{F}_2\) and \(3 \leqslant n \leqslant 5\) the following automorphisms of the vector K-space E are not field automorphisms of E:
\end{enumerate}

\[
\text { for } \quad n = 3, \tau : (\alpha_ {1}, \alpha_ {2}, \alpha_ {3}) \mapsto (\alpha_ {1}, \alpha_ {3}, \alpha_ {2})
\]

\[
\text { for } \quad n = 4, \tau : (\alpha_ {1}, \alpha_ {2}, \alpha_ {3}, \alpha_ {4}) \mapsto (\alpha_ {1}, \alpha_ {4}, \alpha_ {3}, \alpha_ {2})
\]

\[
\text { for } \quad n = 5, \tau : (\alpha_ {1}, \alpha_ {2}, \alpha_ {3}, \alpha_ {4}, \alpha_ {5}) \mapsto (\alpha_ {1}, \alpha_ {5}, \alpha_ {4}, \alpha_ {3}, \alpha_ {2})
\]

but are such that \(\tau(x)\) is conjugate to \(x\) for all \(x \in E\) .

\begin{enumerate}
\def\labelenumi{\arabic{enumi})}
\setcounter{enumi}{7}
\tightlist
\item
  Let \(\mathbf{K}\) be a finite field of \(q\) elements and
\end{enumerate}

\[
\mathrm{P} (\mathbf {X}) = a _ {0} + a _ {1} \mathbf {X} + \dots + a _ {q - 2} \mathbf {X} ^ {q - 2}
\]

a polynomial of K{[}X{]} of degree q - 2. Let r be the rank of the matrix

\[
A = \left( \begin{array}{c c c c} a _ {0} & a _ {1} & \dots & a _ {q - 2} \\ a _ {1} & a _ {2} & \dots & a _ {0} \\ \dots & \dots & \dots & \dots \\ a _ {q - 2} & a _ {0} & \dots & a _ {q - 3} \end{array} \right).
\]

Show that the number of distinct roots of P in K, other than 0, is q - 1 - r. (Consider the product VA, where V is the Vandermonde matrix of elements \(x_{1}, x_{2}, \ldots, x_{q-1}\) of \(K^{*}\) .)

\begin{enumerate}
\def\labelenumi{\arabic{enumi})}
\setcounter{enumi}{8}
\tightlist
\item
  Let \(\mathbf{K}\) be a finite field of \(q\) elements and let \(r = q^n\) for an integer \(n \geqslant 1\) . For every polynomial \(\mathrm{P}(\mathbf{X}) = \sum_{j=0}^{m} a_j \mathbf{X}^j\) of \(\mathbf{K}[\mathbf{X}]\) , put
\end{enumerate}

\[
\hat {\mathrm{P}} (\mathrm{X}) = \sum_ {j} a _ {j} \mathrm{X} ^ {(r ^ {j} - 1) / (r - 1)}
\]

and

\[
\tilde {\mathrm{P}} (\mathrm{X}) = \mathrm{X} \hat {\mathrm{P}} \left(\mathrm{X} ^ {r - 1}\right) = \sum_ {j} a _ {j} \mathrm{X} ^ {r ^ {j}}.
\]

\begin{enumerate}
\def\labelenumi{\alph{enumi})}
\item
  Show that if P, Q are two polynomials in K{[}X{]} and if R = PQ, then we have \(\tilde{\mathbf{R}}(\mathbf{X}) = \tilde{\mathbf{P}}(\tilde{\mathbf{Q}}(\mathbf{X}))\) . Deduce that for two polynomials F, G of K{[}X{]} to be such that F divides G it is necessary and sufficient that \(\hat{F}\) should divide \(\hat{G}\) . (To see that the condition is sufficient, consider the Euclidean division of G by F.)
\item
  Let F be an irreducible polynomial in K{[}X{]} and let G be any polynomial in K{[}X{]}. If there exists an integer \(d \geqslant 1\) such that \(\hat{\mathrm{F}}(\mathbf{X}^{d})\) and \(\hat{\mathrm{G}}(\mathbf{X}^{d})\) have a common root, show that F divides G. (If H is the GCD of \(\hat{\mathrm{F}}(\mathbf{X}^{d})\) and \(\hat{\mathrm{G}}(\mathbf{X}^{d})\) consider the set of polynomials \(P \in K[X]\) such that H divides \(\hat{\mathrm{P}}(\mathbf{X}^{d})\) .)
\item
  Every polynomial \(\mathbf{F} \neq 0\) of \(\mathbf{K}[\mathbf{X}]\) divides a polynomial of the form \(\mathbf{X}^{\mathbf{N}} - 1\) . Show that if \(r - 1 = de\) and if \(\alpha\) is a root of \(\hat{\mathbf{F}}(\mathbf{X}^{d})\) , then \(\alpha\) belongs to the extension \(\mathbf{K}\) of degree \(n\mathbf{N}\) and deduce that the degree of every irreducible factor of \(\hat{\mathbf{F}}(\mathbf{X}^d)\) in \(\mathbf{K}[\mathbf{X}]\) divides \(n\mathbf{N}\) (use \(a\) )). If further \(e\) and \(d\mathbf{N}\) are relatively prime and if \(m\) is the degree of \(\alpha\) over the extension of \(\mathbf{K}\) of degree \(n\) , show that \(m\) divides \(\mathbf{N}\) (use the fact that \((r^{\mathbf{N}} - 1)/e \equiv d\mathbf{N} (\text{mod } e)\) ).
\end{enumerate}

\(d\) ) With the notation as in \(c\) ), suppose that \(\mathbf{F}\) is irreducible and that \(\mathbf{N}\) is the least of the integers \(h\) such that \(\mathbf{F}\) divides \(\mathbf{X}^h - 1\) . If \(e\) and \(d\mathbf{N}\) are relatively prime and if every prime factor of \(n\) divides \(\mathbf{N}\) , show that every irreducible factor of \(\hat{\mathbf{F}}(\mathbf{X}^d)\) is of degree \(n\mathbf{N}\) . (Note that in virtue of \(b\) ) we then have \(m = \mathbf{N}\) .)

\begin{enumerate}
\def\labelenumi{\alph{enumi})}
\setcounter{enumi}{4}
\tightlist
\item
  Deduce in particular from \(d\) ) that if \(\mathrm{P}(\mathbf{X}) = \sum_{j} a_j \mathbf{X}^j\) is an irreducible polynomial in
\end{enumerate}

K{[}X{]} and if N is the least of the integers h such that \(P(X)\) divides \(X^{h}-1\) , then every irreducible factor of \(\sum a_{j}X^{q^{j}-1}\) is of degree N.

\begin{enumerate}
\def\labelenumi{\arabic{enumi})}
\setcounter{enumi}{9}
\tightlist
\item
  Let \(t, n, k\) be three integers \(\geqslant 1\) such that \(t \leqslant n\) and \(n\) divides \(2^k - 1\) . Let \(S\) be the smallest set of integers such that \((1, t) \subset S \subset (1, n)\) and the condition \(\langle i \in S, 2i \equiv j (\text{mod } n) \text{ and } 1 \leqslant j \leqslant n \rangle\) implies \(j \in S\) . Put \(m = \text{Card } S\) . Let \(\xi\) be an element of order \(n\) in \(F_2^k\) and put \(P(X) = \prod_{j \in S} (X - \xi^j)\) .
\end{enumerate}

\begin{enumerate}
\def\labelenumi{\alph{enumi})}
\item
  Prove that \(\mathbf{P}(\mathbf{X})\in \mathbf{F}_2[\mathbf{X}]\) .
\item
  Show that if \(\mathbf{R}(\mathbf{X})\in \mathbf{F}_2[\mathbf{X}]\) is of degree \(< n\) and is a non-zero multiple of \(\mathrm{P}(\mathbf{X})\) , then there are strictly more than \(t\) non-zero coefficients.
\item
  Deduce the construction of a mapping \(\varphi: \mathbf{F}_2^n \to \mathbf{F}_2^m\) such that if \(u, v \in \mathbf{F}_2^n\) are distinct and have at most \(t\) distinct coordinates then \(\varphi(u) \equiv \varphi(v)\) .
\item
  Calculate S, m and P when t = 6, n = 15, k = 4.
\end{enumerate}

\begin{enumerate}
\def\labelenumi{\arabic{enumi})}
\setcounter{enumi}{10}
\tightlist
\item
  Let \(f\) be a monic polynomial in \(\mathbf{Z}[\mathbf{X}]\) , irreducible in \(\mathbf{Q}[\mathbf{X}]\) and let \(p\) be a prime number such that the canonical image \(\vec{f}\) of \(f\) in \(\mathbf{F}_p[\mathbf{X}]\) is a polynomial without multiple roots. If \(\Gamma\) (resp. \(\vec{\Gamma}\) ) is the Galois group of the polynomial \(f\) (resp. \(\vec{f}\) ) over \(\mathbf{Q}\) (resp. \(\mathbf{F}_p\) ), define an isomorphism of a subgroup of \(\Gamma\) onto \(\vec{\Gamma}\) . Deduce that if \(\vec{f} = g_1g_2\ldots g_r\) where the \(g_j\) are irreducible in \(\mathbf{F}_p[\mathbf{X}]\) and if \(g_j\) is of degree \(n_j\) (so that \(f\) has degree \(n = n_1 + n_2 + \dots + n_r\) ), then the group \(\Gamma\) , qua subgroup of \(\mathfrak{S}_n\) contains a product \(\sigma_1\sigma_2\ldots\sigma_r\) of cycles with disjoint supports, \(\sigma_j\) being of order \(n_j\) for \(1\leqslant j\leqslant r\) (cf.~I, p.~62).
\end{enumerate}

Generalize to the case where f is not monic, but where its leading coefficient is not divisible by p.

\begin{enumerate}
\def\labelenumi{\arabic{enumi})}
\setcounter{enumi}{11}
\item
  Let \(f\) be a polynomial of \(\mathbf{Z}[\mathbf{X}]\) of degree \(n\) and let \(p_1, p_2, p_3\) be three distinct prime numbers not dividing the leading coefficient of \(f\) ; let \(f_1, f_2, f_3\) be the canonical images of \(f\) in \(\mathbf{F}_{p_1}[\mathbf{X}], \mathbf{F}_{p_2}[\mathbf{X}], \mathbf{F}_{p_3}[\mathbf{X}]\) respectively. Suppose that \(f_1\) is the product of a linear factor and an irreducible factor of degree \(n - 1\) , \(f_2\) is the product of an irreducible factor of the second degree and irreducible factors of odd degrees, and finally \(f_3\) is irreducible. Show that under these conditions \(f\) is irreducible in \(\mathbf{Q}[\mathbf{X}]\) and that the Galois group of the polynomial \(f\) over \(\mathbf{Q}\) is isomorphic to the symmetric group \(\mathfrak{S}_n\) (« Dedekind's criterion »; use Ex. 11 as well as I, p.~63, Prop. 9).
\item
  Let \(p_1, p_2, \ldots, p_m\) be distinct odd prime numbers and \(P = p_1p_2 \ldots p_m\) ( \(m \geqslant 3\) ).
\end{enumerate}

\begin{enumerate}
\def\labelenumi{\alph{enumi})}
\tightlist
\item
  Given an integer \(n \geqslant 1\) , let \(E\) be the set of polynomials in \(\mathbf{Z}[\mathbf{X}]\) of degree \(\leqslant n\) , whose coefficients lie in the interval \([0, P[\) of \(\mathbf{N}\); we thus have \(\text{Card}(E) = P^{n+1}\) . Let \(k\) be a number such that \(0 < k < 1/3n\) ; show that among the polynomials \(f \in E\) there are at least \(k^m\mathbf{P}^{n + 1}\) such that for every index \(j\) such that \(1 \leqslant j \leqslant m\) , the canonical image \(f_j\) of \(f\) in \(\mathbf{F}_{p_j}[\mathbf{X}]\) is of degree \(n\) and irreducible (use Ex. 1, b)). Suppose moreover that \(k < \frac{1}{18(n - 2)}\) and \(n > 2\) ; arguing in the same way and using Dedekind's criterion (Ex. 12)
\end{enumerate}

show that there exists a set of polynomials \(f \in \mathbf{E}\) , of cardinal at least equal to \((1 - 3(1 - k)^m) \mathbf{P}^{n+1}\) which are of degree \(n\) , irreducible and whose Galois group over \(\mathbf{Q}\) is isomorphic to \(\mathfrak{S}_n\) .

\begin{enumerate}
\def\labelenumi{\alph{enumi})}
\setcounter{enumi}{1}
\tightlist
\item
  For every integer N let \(L_{n,N}\) be the set of polynomials of Z{[}X{]} of degree \(\leqslant n\) , whose coefficients belong to the interval \((-N, N)\) ; we thus have \(\text{Card}(L_{n,N}) = (2N + 1)^{n+1}\) . The integer n being fixed, show that for every \(\varepsilon\) such that \(0 < \varepsilon < 1\) , there exists an integer \(N_0\) such that for \(N \geqslant N_0\) , there are among the polynomials \(f \in L_{n,N}\) at least \((1 - \varepsilon)(2N + 1)^{n+1}\) which are irreducible of degree n, and whose Galois group over Q is isomorphic to \(S_n\) (use a)).
\end{enumerate}

* 14) Let K be a finite field, commutative or not, let Z be its centre, q the number of elements of Z and n the rank {[}K:Z{]}.

\begin{enumerate}
\def\labelenumi{\alph{enumi})}
\item
  If E is a subfield of K containing Z, show that \([E:Z]\) divides n.
\item
  Let \(x \in K\) be an element not belonging to Z; show that the number of distinct conjugates \(yxy^{-1}\) of x in \(K^{*}\) is of the form \((q^{n}-1)/(q^{d}-1)\) , where d divides n and is distinct from n (consider in K the set of elements permutable with x, and use a)).
\item
  Deduce from \(b\) ) that \(q - 1\) is divisible by the integer \(\Phi_n(q)\) (partition the group \(\mathbf{K}^*\) into conjugacy classes and use the relation (6) of \(V\) , p.~82).
\item
  Show that if \(n > 1\) , then \(\Phi_n(q) > (q - 1)^{\varphi(n)}\) (decompose \(\Phi_n(X)\) in the field \(C\) of complex numbers). Deduce that we must have \(K = Z\) , in other words, that \(K\) is necessarily commutative (Wedderburn's theorem). *
\end{enumerate}

\BourbakiExerciseGroup

\begin{enumerate}
\def\labelenumi{\arabic{enumi})}
\tightlist
\item
  Let K be a field of characteristic p \textgreater{} 0 and E an extension of K; E is a p-radical extension of height \(\leqslant 1\) of \(\mathrm{K}(\mathrm{E}^{p})\) . The cardinal of a p-basis of E over \(\mathrm{K}(\mathrm{E}^{p})\) (cf.~V, p.~103, Th. 2) is called the degree of imperfection of E over K. If \(\mathrm{K}(\mathrm{E}^{p}) = \mathrm{E}\) , E is said to be relatively perfect over K; if E is perfect, it is relatively perfect over each of its subfields. The (absolute) degree of imperfection of K is the degree of imperfection of K over its prime subfield \(F_{p}\) , or also the cardinal of a (absolute) p-basis of K.
\end{enumerate}

\begin{enumerate}
\def\labelenumi{\alph{enumi})}
\item
  If B is a \(p\) -basis of E over \(\mathbf{K}(\mathbf{E}^p)\) , show that for every integer \(k > 0\) we have \(\mathbf{E} = \mathbf{K}(\mathbf{E}^{p^k})(\mathbf{B})\) .
\item
  Suppose that \(\mathbf{E} \subset \mathbf{K}^{p^{-n}}\) for an integer \(n > 0\) . Show that for the degree \([\mathbf{E}:\mathbf{K}]\) to be finite it is necessary and sufficient that the degree of imperfection \(m_0\) of \(\mathbf{E}\) over \(\mathbf{K}\) should be finite; then \(m_0\) is the least cardinal of the sets \(\mathbf{S}\) such that \(\mathbf{E} = \mathbf{K}(\mathbf{S})\) . Show that if \(m_k\) is the degree of imperfection of \(\mathbf{K}(\mathrm{E}^{p^k})\) over \(\mathbf{K}\) , then \(m_{k + 1} \leqslant m_k\) for all \(k\) , and if \(f = \sum_{k} m_k\) , then \([\mathbf{E}:\mathbf{K}] = p^f\) .
\end{enumerate}

\begin{enumerate}
\def\labelenumi{\arabic{enumi})}
\setcounter{enumi}{1}
\tightlist
\item
  \begin{enumerate}
  \def\labelenumii{\alph{enumii})}
  \tightlist
  \item
    Let E be an extension of a field K of characteristic p \textgreater{} 0 and F an extension of E. Show that if B is a p-basis of E over \(\mathrm{K}(\mathrm{E}^{p})\) and C a p-basis of F over \(\mathrm{E}(\mathrm{F}^{p})\) , then there exists a p-basis of F over \(\mathrm{K}(\mathrm{F}^{p})\) contained in \(B \cup C\) .
  \end{enumerate}
\end{enumerate}

\begin{enumerate}
\def\labelenumi{\alph{enumi})}
\setcounter{enumi}{1}
\tightlist
\item
  Suppose that B is finite and that F is an algebraic extension of K of finite degree. Show that \([E:K(E^{p})]\geqslant[F:K(F^{p})]\) . If moreover F is separable over E, show that \([E:K(E^{p})]=[F:K(F^{p})]\) .
\end{enumerate}

\begin{enumerate}
\def\labelenumi{\arabic{enumi})}
\setcounter{enumi}{2}
\item
  Let K be an imperfect field and E an algebraic extension of K of finite degree. For an element \(x \in E\) to exist such that \(\mathrm{E} = \mathrm{K}(x)\) it is necessary and sufficient that the degree of imperfection of E over K (Ex. 1) should be 0 or 1. (To see that the condition is sufficient, note that if \(E_{s}\) is the relative separable closure of K in E, we have \(\mathrm{E} = \mathrm{E}_{s}(\alpha)\) and \(\mathrm{E}_{s} = \mathrm{K}(\beta)\) and if \(\alpha \notin E_{s}\) , consider the conjugates of \(\alpha + c\beta\) over K, for \(c \in K\) .)
\item
  Let K be a field of characteristic p \textgreater{} 0 and E an algebraic extension of K of finite degree. If r \textgreater{} 0 is the degree of imperfection of E over K (Ex. 1), show that r is the least cardinal of sets S such that \(\mathrm{E} = \mathrm{K}(\mathrm{S})\) . (To see that E can be generated by r of its elements note that if \(E_{s}\) is the relative separable closure of K in E, then there exist r elements \(a_{i}\) ( \(1 \leqslant i \leqslant r\) ) of E such that \(\mathrm{E} = \mathrm{E}_{s}(a_{1}, \ldots, a_{r})\) , and use Ex. 3.)
\item
  Let K be a field and E an algebraic extension of K of finite degree.
\end{enumerate}

\begin{enumerate}
\def\labelenumi{\alph{enumi})}
\item
  Suppose that K is of characteristic p \textgreater{} 0. Show that if the degree of imperfection of E over K (Ex. 1) is \textgreater{} 1, there exist an infinity of distinct fields F such that \(K \subset F \subset E\) (reduce to the case where \(\mathrm{K}(\mathrm{E}^{p}) = \mathrm{K}\) , and consider the fields \(\mathrm{K}(a + \lambda b)\) , where the set \(\{a, b\}\) is p-free over K and \(\lambda \in K\) ).
\item
  Conclude that for there to exist only a finite number of fields F such that \(K \subset F \subset E\) it is necessary and sufficient that the degree of imperfection of E over K be \(\leqslant 1\) .
\end{enumerate}

\begin{enumerate}
\def\labelenumi{\arabic{enumi})}
\setcounter{enumi}{5}
\tightlist
\item
  Let K be a field of characteristic p \textgreater{} 0, E an algebraic extension of K of finite degree and N the quasi-Galois extension of K generated by E.
\end{enumerate}

\begin{enumerate}
\def\labelenumi{\alph{enumi})}
\item
  Let \(E_{r}\) and \(E_{s}\) (resp. \(N_{r}\) and \(N_{s}\) ) be the largest p-radical extension and the largest separable extension of K contained in E (resp. in N). Show that for \(\mathrm{E} = \mathrm{E}_{r}(\mathrm{E}_{s})\) (in which case E is isomorphic to \(E_{s} \otimes_{K} E_{r}\) ) it is necessary and sufficient that \(E_{r} = N_{r}\) .
\item
  If the degree of imperfection of E over K is 1, then a necessary and sufficient condition for \(\mathrm{E} = \mathrm{E}_{r}(\mathrm{E}_{s})\) is that the degree of imperfection of N over K be 1. Deduce that if N is a quasi-Galois extension of K, whose degree of imperfection over K is equal to 1, then \(\mathrm{E} = \mathrm{E}_{r}(\mathrm{E}_{s})\) for every subextension E of N (cf.~Ex. 2, b)).
\item
  Conversely if N is a quasi-Galois extension of K of finite degree whose degree of imperfection over K is \textgreater1, show that if \(N_{r} \neq N\) , then there exists \(t \in N\) such that for the extension \(\mathrm{E} = \mathrm{K}(t)\) of K we have \(\mathrm{E} \neq \mathrm{E}_{r}(\mathrm{E}_{s})\) (if a, b are two elements of a p-basis of N over \(\mathrm{K}(\mathrm{N}^{p})\) and if \(x \in N\) is separable over K and \(x \notin K\) , take \(t = a + bx\) ).
\end{enumerate}

\BourbakiExerciseGroup

\begin{enumerate}
\def\labelenumi{\arabic{enumi})}
\item
  Let E be an extension of a field K and let B be a transcendence basis of E over K. Show that E is equipotent to \(K \times B\) if one of the sets K, B is infinite, and countable otherwise. * Deduce in particular that every transcendence basis of the field R of real numbers over the field Q of rational numbers has the power of the continuum. *
\item
  Let K be the field \(\mathbf{Q}(\mathbf{X})\) of rational fractions in one indeterminate over the field Q of rational numbers. Show that in the ring K{[}Y{]} the polynomial \(Y^{2} + X^{2} + 1\) is irreducible and if E is the extension of K generated by a root of this polynomial (in an algebraic closure of
\end{enumerate}

K), then Q is relatively algebraically closed in E but E is not a pure extension of Q. (To see that there is in E no element \(a \notin Q\) algebraic over Q note that the existence of such an element would imply the relation \(\mathrm{E} = \mathbf{Q}(a)(\mathbf{X})\) , and observe that \(-(\mathbf{X}^{2} + 1)\) is not a square in \(\Omega(\mathbf{X})\) , where \(\Omega\) is an algebraic closure of Q.) Show that if i is a root of \(X^{2} + 1\) , then \(\mathrm{E}(i)\) is a pure transcendental extension of \(\mathbf{Q}(i)\) .

* 3) Let K be the field C(X) of rational fractions in one indeterminate over the (algebraically closed) field C of complex numbers. Show that in the ring K{[}Y{]}, the polynomial \(Y^{3} + X^{3} + 1\) is irreducible; let E be the extension of K generated by a root of this polynomial (in an algebraic closure of K). Show that E is not a pure extension of C, even though C is algebraically closed. (Show that it is impossible for a relation \(u^{3} + v^{3} + w^{3} = 0\) to exist between three polynomials u, v, w of C{[}X{]}, pairwise relatively prime and not all three constant. Argue by contradiction: if r is the largest of the degrees of u, v, w and if for example \(\deg(w) = r\) , deduce from the relation

\[
w ^ {3} = - (u + v) (u + j v) (u + j ^ {2} v)
\]

where \(j\) is a primitive cube root of unity, the existence of three polynomials \(u_1\) , \(v_1\) , \(w_1\) in \(\mathbf{C}[\mathbf{X}]\) , pairwise relatively prime, not all three constant, of degrees \(< r\) and such that \(u_1^3 + v_1^3 + w_1^3 = 0\) .) *

\begin{enumerate}
\def\labelenumi{\arabic{enumi})}
\setcounter{enumi}{3}
\tightlist
\item
  Let \(\Omega\) be an algebraically closed extension of a field \(\mathbf{K}\) , of infinite transcendence degree over \(\mathbf{K}\) .
\end{enumerate}

\begin{enumerate}
\def\labelenumi{\alph{enumi})}
\tightlist
\item
  Show that there exist an infinity of K-endomorphisms of \(\Omega\) having as images subfields of \(\Omega\) distinct from \(\Omega\) , and with respect to which \(\Omega\) may have a transcendence degree equal to any cardinal up to at most tr. \(\deg_{\mathbf{K}}\Omega\) . * In particular, there exist an infinity of distinct Q-isomorphisms of the field C of complex numbers onto subfields of C distinct from C. * b) Show that for every subextension E of \(\Omega\) such that tr. \(\deg_{\mathbf{K}}E < \text{tr} \cdot \deg_{\mathbf{K}}\Omega\) , every K-isomorphism of E onto a subfield of \(\Omega\) may be extended to a K-automorphism of \(\Omega\) . Give examples where tr. \(\deg_{\mathbf{K}}E = \text{tr} \cdot \deg_{\mathbf{K}}\Omega\) and where there exists a K-isomorphism of E onto a subfield of \(\Omega\) which cannot be extended to a K-isomorphism of any extension F of E, contained in \(\Omega\) and distinct from E, onto a subfield of \(\Omega\) .
\end{enumerate}

\begin{enumerate}
\def\labelenumi{\arabic{enumi})}
\setcounter{enumi}{4}
\item
  Let \(\Omega\) be an algebraically closed extension of a field \(\mathbf{K}\) . Show that for every subextension \(\mathbf{E}\) of \(\Omega\) , which is transcendental over \(\mathbf{K}\) , there exist an infinity of \(\mathbf{K}\) -isomorphisms of \(\mathbf{E}\) onto a subfield of \(\Omega\) (if \(x \in \mathbf{E}\) is transcendental over \(\mathbf{K}\) , consider a subextension \(\mathbf{F}\) of \(\mathbf{E}\) such that \(x\) is transcendental over \(\mathbf{F}\) and \(\mathbf{E}\) algebraic over \(\mathbf{F}(x)\) , and show that there exist an infinity of \(\mathbf{F}\) -isomorphisms of \(\mathbf{E}\) onto a subfield of \(\Omega\) ).
\item
  Let \(\Omega\) be an algebraically closed extension of a field \(K\) , and \(N\) a subextension of \(\Omega\) . Show that if \(N\) is quasi-Galois over \(K\) and if \(E\) and \(E'\) are two conjugate extensions of \(K\) contained in \(\Omega\) , then \(N(E)\) and \(N(E')\) are conjugate over \(K\) . Conversely, if \(N\) has this property and if \(\operatorname{tr} \cdot \deg_{K} N\) is finite and \(< \operatorname{tr} \cdot \deg_{K} \Omega\) , then it is necessarily algebraic and quasi-Galois over \(K\) .
\item
  Let E be an extension of a field K and \(\operatorname{Aut}_{\mathbf{K}}(\mathbf{E})\) the group of all K-automorphisms of E. a) For every finitely generated subextension L of E over K, \(\operatorname{Aut}_{\mathbf{L}}(\mathbf{E})\) is a subgroup of \(\operatorname{Aut}_{\mathbf{K}}(\mathbf{E})\) . Show that when L ranges over the set of subextensions of E finitely generated over K, the groups \(\operatorname{Aut}_{\mathbf{L}}(\mathbf{E})\) form a fundamental system of neighbourhoods of the neutral element of \(\mathrm{Aut}_{\mathbf{K}}(\mathbf{E})\) for a topology compatible with the group structure of \(\mathrm{Aut}_{\mathbf{K}}(\mathbf{E})\) and for which \(\mathrm{Aut}_{\mathbf{K}}(\mathbf{E})\) is a separated and totally disconnected group.
\end{enumerate}

\begin{enumerate}
\def\labelenumi{\alph{enumi})}
\setcounter{enumi}{1}
\tightlist
\item
  When E is equipped with the discrete topology, the topology of \(\operatorname{Aut}_{\mathbf{K}}(\mathbf{E})\) is the coarsest in which the mapping \((u,x)\to u(x)\) of \(\operatorname{Aut}_{\mathbf{K}}(\mathbf{E})\times\mathbf{E}\) into E is continuous.
\end{enumerate}

* c) The group \(\operatorname{Aut}_{\mathbf{Q}}(\mathbf{R})\) is reduced to the neutral element (V, p.~153, Ex. 2 of § 9). * \(d)\) Let \(\mathbf{K}_0 \supseteq \mathbf{K}\) be the subextension of E consisting of the elements invariant under all K-automorphisms of E. Show that for E to be algebraic over \(\mathbf{K}_0\) it is necessary and sufficient that \(\operatorname{Aut}_{\mathbf{K}}(\mathbf{E})\) should be compact.

\begin{enumerate}
\def\labelenumi{\alph{enumi})}
\setcounter{enumi}{4}
\tightlist
\item
  Assume that K is infinite. If \(E = K(X)\) , the field of fractions in one indeterminate over K, show that \(\operatorname{Aut}_{K}(E)\) is an infinite discrete group (V, p.~148, Ex. 11). If K is of characteristic 0, the subgroup \(\Gamma\) of \(\operatorname{Aut}_{K}(E)\) generated by the automorphism \(\sigma : X \mapsto X + 1\) is closed in \(\operatorname{Aut}_{K}(E)\) and distinct from the latter, but has the same field of invariants.
\end{enumerate}

\begin{enumerate}
\def\labelenumi{\arabic{enumi})}
\setcounter{enumi}{7}
\tightlist
\item
  Let \(E\) be an extension of a field \(K\) .
\end{enumerate}

\begin{enumerate}
\def\labelenumi{\alph{enumi})}
\item
  Let \(x\) be an element of \(E - K\) . Show that there exists a pure extension \(L \subset E\) of \(K\) such that \(E\) is algebraic over \(L\) and \(x \notin L\) (consider the set of subextensions \(L\) of \(E\) which are pure transcendental over \(K\) and such that \(x \notin L\) ).
\item
  Let \(x \in E\) be a transcendental element over \(K\) and \(p\) an integer \(\geqslant 2\) . Show that there exists a pure transcendental extension \(L \subset E\) of \(K\) such that \(E\) is algebraic over \(L\) , \(x \notin L\) and \(x^p \in L\) (same method).
\item
  Let F be a subextension of E and let \(x \in E\) be an algebraic element over F and \(P \in F[X]\) its minimal polynomial. Show that there exists a pure extension \(L \subset E\) of K such that E is algebraic over L and P is irreducible over the field \(F(L)\) .
\end{enumerate}

\begin{enumerate}
\def\labelenumi{\arabic{enumi})}
\setcounter{enumi}{8}
\tightlist
\item
  An extension E of a field K is called a Dedekind extension of K if for every subextension F of E, F is identical with the subfield of elements of E invariant under the group \(\mathrm{Aut}_{\mathrm{F}}(\mathrm{E})\) of all F-automorphisms of E.
\end{enumerate}

\begin{enumerate}
\def\labelenumi{\alph{enumi})}
\item
  For an algebraic extension \(\mathbf{E}\) of \(\mathbf{K}\) to be Dedekind over \(\mathbf{K}\) it is necessary and sufficient that \(\mathbf{E}\) should be a Galois extension of \(\mathbf{K}\) .
\item
  Suppose that K is of characteristic p \textgreater{} 0. Show that a Dedekind extension of K is necessarily algebraic over K (hence Galois). (Use Ex. 8, b).)
\item
  If E is a Dedekind extension of K and L a pure transcendental subextension of E, distinct from K and such that E is an algebraic extension of L, show that E has infinite degree over L.
\end{enumerate}

\begin{enumerate}
\def\labelenumi{\arabic{enumi})}
\setcounter{enumi}{9}
\item
  Let E be an extension of a field K such that the mapping \(F \mapsto \operatorname{Aut}_{F}(E)\) is a bijection of the set of subextensions of E onto the set of subgroups of \(\operatorname{Aut}_{K}(E)\) . Show that E is then a Galois extension of K of finite degree. (Using Ex. 9, c), show first that E is necessarily algebraic over K, then use Ex. 9, a).
\item
  Let \(\mathbf{E}\) be an extension of a field \(\mathbf{K}\) and \(\Omega\) an algebraically closed extension of \(\mathbf{E}\) . Show that the following conditions are equivalent:
\end{enumerate}

α) E is a p-radical extension of K;

\(\beta)\) for every extension \(F\) of \(K\) , \(E \otimes_{K} F\) has only a single prime ideal;

γ) for every extension F of K, any two composite extensions of E and F are isomorphic;

δ) \(E \otimes_{K} \Omega\) has only a single prime ideal ;

ε) any two composite extensions of E and Ω are isomorphic.

\begin{enumerate}
\def\labelenumi{\arabic{enumi})}
\setcounter{enumi}{11}
\tightlist
\item
  \begin{enumerate}
  \def\labelenumii{\alph{enumii})}
  \tightlist
  \item
    Let K be a field, \(\Omega\) an algebraically closed extension of K, \(E \subset \Omega\) and F two extensions of K; assume that \(\operatorname{tr} \cdot \deg_E \Omega \geqslant \operatorname{tr} \cdot \deg_K F\) . Show that every composite extension of E and F is isomorphic to a composite of the form \((L_w, 1, w)\) , where \(w\) is a K-isomorphism of F onto a subfield of \(\Omega\) and \(L_w\) is the subfield of \(\Omega\) generated by E and \(w(F)\) .
  \end{enumerate}
\end{enumerate}

\begin{enumerate}
\def\labelenumi{\alph{enumi})}
\setcounter{enumi}{1}
\tightlist
\item
  Deduce from a) that if F is algebraic of finite degree n over K, then there are at most n composite extensions of E and F that are pairwise non-isomorphic.
\end{enumerate}

\begin{enumerate}
\def\labelenumi{\arabic{enumi})}
\setcounter{enumi}{12}
\item
  In an algebraic closure \(\Omega\) of \(\mathbf{Q}(\mathbf{X})\) consider two pure transcendental extensions \(\mathrm{E} = \mathbf{Q}(\mathbf{X})\) and \(\mathrm{F} = \mathbf{Q}(\mathbf{X} + i)\) (where \(i^2 = -1\) ) of the field \(\mathbf{Q}\) . Show that \(\mathrm{E} \cap \mathrm{F} = \mathbf{Q}\) but that \(\mathrm{E}\) and \(\mathrm{F}\) are not algebraically disjoint over \(\mathbf{Q}\) .
\item
  Let E, F, G be three extensions of a field K contained in an extension \(\Omega\) of K and such that \(F \subset G\) . Show that for E and G to be algebraically disjoint over K it is necessary and sufficient that E and F should be algebraically disjoint over K and \(\mathrm{E}(\mathrm{F})\) and G algebraically disjoint over F.
\item
  \begin{enumerate}
  \def\labelenumii{\alph{enumii})}
  \tightlist
  \item
    Let K be a field and L a subfield of K such that K is a finitely generated L-algebra, in other words \(K = L[a_{1}, a_{2}, \ldots, a_{n}]\) for elements in K. Show that the \(a_{j}\) are algebraic over L. (Argue by contradiction: if \(a_{1}, \ldots, a_{m}\) ( \(m \geqslant 1\) ) form a maximal algebraically free subfamily of \((a_{j})_{1 \leqslant j \leqslant n}\) , observe that in the ring \(A = L[a_{1}, \ldots, a_{m}]\) the intersection of all non-zero ideals is reduced to zero, and use Ex. 5 of V, p.~147.)
  \end{enumerate}
\end{enumerate}

\begin{enumerate}
\def\labelenumi{\alph{enumi})}
\setcounter{enumi}{1}
\tightlist
\item
  Deduce from a) that in the polynomial algebra \(K[X_{1}, \ldots, X_{n}]\) every maximal ideal is of finite codimension over K. Conclude that if A is a commutative finitely generated algebra over K, then every subfield of A containing K is of finite degree over K.
\end{enumerate}

\begin{enumerate}
\def\labelenumi{\arabic{enumi})}
\setcounter{enumi}{15}
\tightlist
\item
  \begin{enumerate}
  \def\labelenumii{\alph{enumii})}
  \tightlist
  \item
    Let K be an algebraically closed field, \(\Omega\) an algebraically closed extension of K, E a subextension of \(\Omega\) and \(x\) an element of \(\Omega\) transcendental over E. Let \(\overline{\mathbf{K}(x)}\) be the relative algebraic closure of \(\mathbf{K}(x)\) in \(\Omega\) and let \(\mathbf{M} = \operatorname{E}(\overline{\mathbf{K}(x)})\) ; show that M is an algebraic extension of infinite degree of \(\operatorname{E}(x)\) . (Observe that for every integer \(m > 1\) the polynomial \(\mathbf{X}^m - x\) is irreducible in \(\operatorname{E}(x)[\mathbf{X}]\) (V, p.~91, Remark).
  \end{enumerate}
\end{enumerate}

\begin{enumerate}
\def\labelenumi{\alph{enumi})}
\setcounter{enumi}{1}
\tightlist
\item
  Let E be an algebraically closed field and F a finitely generated extension of E. Show that every algebraically closed subfield of F is necessarily contained in E. (Otherwise there would exist a subfield \(L \subset F\) which is algebraically closed and an element \(x \in L\) which is transcendental over E. Apply a), taking for K the relative algebraic closure in E of the prime subfield of E, noting that M should be finitely generated over E.)
\end{enumerate}

\begin{enumerate}
\def\labelenumi{\arabic{enumi})}
\setcounter{enumi}{16}
\item
  Let K be a field, \(\Omega\) an algebraically closed extension of K and x an element of \(\Omega\) transcendental over K. Show that for every prime number q there exists a subextension M of \(\Omega\) such that \(\mathbf{M}(x)\) is a Galois extension of degree q of M. (Consider the element \(y = x^{q}\) if q is not the characteristic of K, \(y = x^{q} - x\) in the opposite case, and a subextension L of \(\Omega\) generated by a transcendence basis of \(\Omega\) over K containing y and besides the q-th roots of unity if \(y = x^{q}\) ; use Ex. 12 of V, p.~163.)
\item
  Let K be a field and \(\mathbf{K}((\mathbf{X}))\) the field of formal power series in one indeterminate over K (IV, p.~38). Show that \(\operatorname{tr}.\deg_{\mathbf{K}}\mathbf{K}((\mathbf{X}))=\operatorname{Card}(\mathbf{K}^{\mathbf{N}})\) . Distinguish two cases:
\end{enumerate}

\begin{enumerate}
\def\labelenumi{\alph{enumi})}
\item
  If \(\operatorname{Card}(\mathbf{K}) < \operatorname{Card}(\mathbf{K}^{\mathbf{N}})\) , note that \(\operatorname{Card}(\mathbf{K}((\mathbf{X}))) = \operatorname{Card}(\mathbf{K}^{\mathbf{N}})\) .
\item
  If \(\operatorname{Card}(K) = \operatorname{Card}(K^{N})\) , let P be the prime subfield of K, S an infinite set of elements of \(\operatorname{K}((X))\) which is algebraically independent over K and T the set of coefficients of all the formal power series belonging to S. Let L be the relative algebraic closure in
\end{enumerate}

K((X)) of the field K(S) and u an element of L; there is an equation \(g(s_{1}, \ldots, s_{m}, u) = 0\) where \(g \neq 0\) is a polynomial in \(K[X_{1}, \ldots, X_{m}, X_{m+1}]\) and \(s_{1}, \ldots, s_{m}\) elements of S. Let A be the set of coefficients of the polynomial g and \(C(u)\) the set of coefficients of the formal power series u; show that the field \(P(T)(A \cup C(u))\) is algebraic over \(P(T \cup A)\) (denoting by \(\Omega\) an algebraic closure of K, show that otherwise there would exist an infinity of formal power series \(v \in \Omega((X))\) satisfying the equation g ( \(s_{1}, \ldots, s_{m}, v) = 0\) ). Using Ex. 1, show that if \(\text{Card}(S) < \text{Card}(K)\) , then the transcendence degree of K over \(P(T)\) is infinite and deduce that in this case L is distinct from \(K((X))\) (if \((t_{n})_{n \geqslant 0}\) is an infinite sequence of elements of K algebraically independent over \(P(T)\) , consider the formal power series \(\sum_{n=0}^{\infty} t_{n} X^{n}\) ).

\begin{enumerate}
\def\labelenumi{\arabic{enumi})}
\setcounter{enumi}{18}
\tightlist
\item
  Let \(\mathbf{E}\) and \(\mathbf{F}\) be two transcendental extensions of a field \(\mathbf{K}\) , linearly disjoint over \(\mathbf{K}\) . Show that \(\mathbf{K}(\mathbf{E} \cup \mathbf{F})\) is distinct from the ring \(\mathbf{C}\) (isomorphic to \(\mathbf{E} \otimes_{\mathbf{K}} \mathbf{F}\) ) generated by \(\mathbf{E} \cup \mathbf{F}\) . (Reduce to the case where the transcendence degrees of \(\mathbf{E}\) and \(\mathbf{F}\) over \(\mathbf{K}\) are equal to 1; if \(x \in \mathbf{E}\) and \(y \in \mathbf{F}\) are transcendental over \(\mathbf{K}\) show that we cannot have \((x + y)^{-1} \in \mathbf{C}\) ; arguing by contradiction, show that in the opposite case, if \(p\) is the characteristic exponent of \(\mathbf{K}\) , there would exist an integer \(r \geqslant 0\) such that \((x + y)^{-p^r}\) belonged to the subring of \(\mathbf{C}\) generated by \(\mathbf{K}(x) \cup \mathbf{K}(y)\) .)
\end{enumerate}

* 20) Let \(n\) be an integer, \(K\) a field of characteristic \(p\) prime to \(n\) , possessing primitive \(n\) -th roots of unity, \(\mu_n \subset K^*\) the subgroup of \(n\) -th roots of unity, \(G\) a finite commutative group annihilated by \(n\) and \(A = K^{(G)}\) the group algebra of \(G\) (III, p.~446).

\begin{enumerate}
\def\labelenumi{\alph{enumi})}
\item
  Show that A is a diagonalizable algebra (V, p.~28). (Decompose G into a direct sum of cyclic groups (VII, p.~22).) Deduce that every A-module is a direct sum of modules which are vector K-spaces of dimension 1.
\item
  Let V be an A-module of finite dimension over K, S the symmetric algebra of the vector K-space V and F the field of fractions of S. The group G operates on V, hence on S and F. Assume that G acts faithfully on V. Show that G operates faithfully on F. Calculate \([F : F^{G}]\) .
\item
  Let \((x_{1},\ldots ,x_{n})\) be a basis of the vector K-space V consisting of eigenvectors for the action of G. Let \(\Gamma\) be the subgroup of \(\mathbf{F}^*\) generated by \(x_{1},\ldots ,x_{n}\) . Show that \(\Gamma\) is a free commutative group and that the sub-K-algebra B of F generated by \(\Gamma\) may be identified with \(\mathbf{K}^{(\Gamma)}\) and is stable under the operation of G.
\item
  For all \(x \in \Gamma\) , the mapping \(\gamma \mapsto \frac{\gamma x}{x}\) is a homomorphism of \(G\) into \(\mu_n\) , giving rise to a homomorphism \(\theta: \Gamma \to \operatorname{Hom}(G, \mu_n)\) . Show that \(\theta\) is surjective. Put \(\Gamma_1 = \operatorname{Ker}(\theta)\) . Using the elementary divisor theorem (VII, p.~24), show that \(B\) is a free \(K^{(\Gamma_1)}\) -module of rank Card(G).
\item
  By passing to the field of fractions deduce from \(d\) and \(a\) that \(\mathbf{F}^{\mathrm{G}}\) is the field of fractions of \(\mathbf{K}^{(\Gamma_1)}\) . Deduce that \(\mathbf{F}^{\mathrm{G}}\) is a pure transcendental extension of \(\mathbf{K}\) . *
\end{enumerate}

\BourbakiExerciseGroup

\begin{enumerate}
\def\labelenumi{\arabic{enumi})}
\tightlist
\item
  Let \(\mathbf{K}\) be a field of characteristic \(p > 0\) , \(\mathbf{E}\) an extension of \(\mathbf{K}\) and \(\mathbf{B}\) a \(p\) -basis of \(\mathbf{E}\) over \(\mathbf{K}(\mathbf{E}^p)\) (V, p.~98).
\end{enumerate}

\begin{enumerate}
\def\labelenumi{\alph{enumi})}
\item
  Assume that \(\mathbf{E} \subset \mathbf{K}^{p^{-1}}\) ; let \(\mathbf{K}_0\) be a subfield of \(\mathbf{K}\) such that \(\mathbf{E}\) is separable over \(\mathbf{K}_0\) . Show that the set \(\mathbf{B}^p\) is a \(p\) -independent subset over \(\mathbf{K}_0(\mathbf{K}^p)\) (observe that if \((a_{\lambda})\) is a basis of the vector \(K_{0}\) -space K, then \((a_{\lambda}^{p})\) is a basis of \(K_{0}(K^{p})\) over \(K_{0}\) . Let C be a subset of K, disjoint from \(B^{p}\) and such that \(B^{p} \cup C\) is a p-basis of K over \(K_{0}(K^{p})\) ; show that \(B \cup C\) is a p-basis of E over \(K_{0}(E^{p})\) .
\item
  Assume that \(E \subset K^{p^{-n}}\) and that E is separable over a subfield \(K_{0}\) of K. Show that if the degree of imperfection of K over \(K_{0}\) is finite (V, p.~170. Ex. 1), then it is equal to the degree of imperfection of E over \(K_{0}\) (use a)).
\end{enumerate}

\begin{enumerate}
\def\labelenumi{\arabic{enumi})}
\setcounter{enumi}{1}
\tightlist
\item
  Let \(\mathbf{K}\) be a field of characteristic \(p > 0\) , \(\mathbf{E}\) an extension of \(\mathbf{K}\) and \(\mathbf{F}\) a separable extension of \(\mathbf{E}\) . Show that any two of the following three properties imply the third:
\end{enumerate}

\begin{enumerate}
\def\labelenumi{\alph{enumi})}
\item
  B is a \(p\) -basis of E over \(\mathbf{K}(\mathbf{E}^p)\) ;
\item
  C is a \(p\) -basis of \(\mathbf{F}\) over \(\operatorname{E}(\mathbf{F}^p)\) ;
\item
  \(\mathbf{B} \subset \mathbf{E}\) , \(\mathbf{B} \cup \mathbf{C}\) is a \(p\) -basis of \(\mathbf{F}\) over \(\mathbf{K}(\mathbf{F}^p)\) and \(\mathbf{B} \cap \mathbf{C} = \emptyset\) . (Use the fact that if \((c_{\mu})\) is a basis of the vector E-space \(\mathbf{F}\) , then \((c_{\mu}^p)\) is a basis of \(\mathbf{K}(\mathbf{E}^p)[\mathbf{F}^p]\) over \(\mathbf{K}(\mathbf{E}^p)\) and of \(\mathbf{E}[\mathbf{F}^p]\) over \(\mathbf{E}\) .)
\end{enumerate}

\begin{enumerate}
\def\labelenumi{\arabic{enumi})}
\setcounter{enumi}{2}
\item
  Let K be a field of characteristic p \textgreater{} 0, E an extension of K, and F a finitely generated extension of E. Show that the degree of imperfection of F over K is at least equal to that of E over K (reduce to the following two cases: \(1^{\circ}\) F is separable over E; \(2^{\circ}\) F = E(x), where \(x^{p} \in E\) (cf.~V, p.~170, Ex. 2)).
\item
  Let \(\mathbf{F}\) be a separable extension of a field \(\mathbf{K}\) of characteristic \(p > 0\) . Show that if \(\mathbf{E}\) is a subextension of \(\mathbf{F}\) , relatively perfect over \(\mathbf{K}\) (V, p.~170, Ex. 1), \(\mathbf{F}\) is separable over \(\mathbf{E}\) .
\item
  Let \(\mathbf{K}\) be a field of characteristic \(p > 0\) . Show that if \(\mathbf{E}\) is an algebraic extension of \(\mathbf{K}\) or a relatively perfect extension of \(\mathbf{K}\) , then \(\mathrm{E}^{p^{-\infty}} = \mathrm{E}(\mathrm{K}^{p^{-\infty}})\) .
\item
  Let \(\mathbf{K}\) be a field of characteristic \(p > 0\) , \(\mathbf{E}\) a separable extension of \(\mathbf{K}\) and \(\mathbf{B}\) a \(p\) -basis of \(\mathbf{E}\) over \(\mathbf{K}(\mathbf{E}^p)\) .
\end{enumerate}

\begin{enumerate}
\def\labelenumi{\alph{enumi})}
\item
  Show that B is algebraically free over K (consider an algebraic relation of least degree between the elements of B and put the degrees of the variables which occur in the form \(kp + h\) with \(0 \leqslant h \leqslant p - 1\) ). Deduce that if tr. \(\deg_{K}E < +\infty\) , then the degree of imperfection of E over K is at most equal to tr. \(\deg_{K}E\) .
\item
  Show that E is separable and relatively perfect over K(B).
\end{enumerate}

\begin{enumerate}
\def\labelenumi{\arabic{enumi})}
\setcounter{enumi}{6}
\item
  Show that a relatively perfect transcendental extension E of a field K of characteristic p \textgreater{} 0 is not finitely generated over K. (In the opposite case show that for every transcendence basis B of E over K, E would be algebraic and separable over K(B).)
\item
  If E is a separable algebraic extension of a field K of characteristic p \textgreater{} 0, then the largest perfect subextension \(E^{p^{\infty}}\) of E (the intersection of the fields \(E^{p^{n}}\) for \(n \geqslant 0\) ) is algebraic over the largest perfect subfield \(K^{p^{\infty}}\) of K. (If \(x \in E^{p^{\infty}}\) , show that for every integer n \textgreater{} 0 we have \(x^{p^{-n}} \in \mathrm{K}(x)\) and deduce that the minimal polynomial of x over \(K^{p^{n}}\) is the same as its minimal polynomial over K.)
\item
  Let \(\mathbf{E}\) be a field, \(\mathbf{G}\) a group of automorphisms of \(\mathbf{E}\) and \(\mathbf{K}\) the subfield of \(\mathbf{E}\) consisting of the elements invariant under \(\mathbf{G}\) .
\end{enumerate}

\begin{enumerate}
\def\labelenumi{\alph{enumi})}
\item
  Show that for every subfield L of E which is stable for the elements of G, L and K are linearly disjoint over L ∩ K. (Consider a linear relation \(\sum_{j}\lambda_{j}x_{j}=0\) between the elements \(x_{j} \in \mathbf{K}\) which are linearly independent over \(\mathbf{L} \cap \mathbf{K}\) , with \(\lambda_{j} \in \mathbf{L}\) not all zero, the number of \(\lambda_{j} \neq 0\) being the least possible.)
\item
  If \(L \cap K\) is the field of invariants of a group of automorphisms of K, show that L is the field of invariants of a group of automorphisms of \(\mathrm{L}(\mathbf{K})\) .
\end{enumerate}

\begin{enumerate}
\def\labelenumi{\arabic{enumi})}
\setcounter{enumi}{9}
\item
  Let E be an extension of a field K, let G be a compact subgroup of the topological group \(\operatorname{Aut}_{\mathbf{K}}(\mathbf{E})\) and let F be the field of invariants of G. Show that E is a Galois extension of F and that G is canonically isomorphic to the topological group \(\operatorname{Aut}_{\mathbf{F}}(\mathbf{E})\) . (Observe that G operates in a continuous fashion on the discrete space E to establish that an element \(x \in E\) has a finite orbit under the action of G.)
\item
  Let \(\mathbf{K}\) be a field of characteristic \(p > 0\) . For every commutative K-algebra \(\mathbf{A}\) , raising to the \(p\) -th power makes the ring \(\mathbf{A}\) into an A-algebra, hence also a K-algebra which is denoted by \(\mathbf{A}^{p^{-1}}\) . Let \(\mathbf{A}\) be a commutative K-algebra. Show that the following conditions are equivalent:
\end{enumerate}

\begin{enumerate}
\def\labelenumi{(\roman{enumi})}
\item
  \(\mathbf{A}\) is a separable K-algebra.
\item
  The unique K-homomorphism \(\Phi: \mathbf{A} \otimes_{\mathbf{K}} \mathbf{K}^{p^{-1}} \to \mathbf{A}^{p^{-1}}\) such that \(\Phi(a \otimes \lambda) = a^p \cdot \lambda (a \in \mathbf{A}, \lambda \in \mathbf{K})\) is injective.
\item
  For every family \((b_{i})_{i\in I}\) of \(\mathbf{K}\) which is linearly free over \(\mathbf{K}^p\) and every family \((a_{i})_{i\in I}\) of \(\mathbf{A}\) the equality \(\sum_{i}a_i^p b_i = 0\) implies that \(a_{i} = 0\) for all \(i\) .
\end{enumerate}

(It may be verified that (ii) \(\Leftrightarrow\) (iii) and that (ii) is a consequence of V, p.~123, Th. 2, e).)

\begin{enumerate}
\def\labelenumi{\arabic{enumi})}
\setcounter{enumi}{11}
\tightlist
\item
  Let K be a field and \((\mathbf{X}_{i})_{i\in I}\) a family of indeterminates. Show that the algebra of formal power series \(\mathbf{K}[[(\mathbf{X}_{i})_{i\in I}]]\) is a separable K-algebra. (Apply Ex. 11; it may also be verified that for every finite extension \(K'\) of K, \(\mathbf{K}[[(\mathbf{X}_{i})_{i\in I}]]\otimes_{\mathbf{K}}K'\) is isomorphic to \(\mathbf{K}'[[(\mathbf{X}_{i})_{i\in I}]]\) and then Th. 2, d) of V, p.~123 applied.) Deduce that the field of fractions of \(\mathbf{K}[[(\mathbf{X}_{i})_{i\in I}]]\) is a separable extension of K (V, p.~120, Prop. 4).
\end{enumerate}

\BourbakiExerciseGroup

\begin{enumerate}
\def\labelenumi{\arabic{enumi})}
\tightlist
\item
  Let \(\mathbf{K}\) be a field of characteristic \(p > 0\) , E a separable extension of \(\mathbf{K}\) , of finite transcendence degree over \(\mathbf{K}\) .
\end{enumerate}

\begin{enumerate}
\def\labelenumi{\alph{enumi})}
\item
  If there exists a transcendence basis \(B_{0}\) of E over K and an integer \(m \geqslant 0\) such that \(\mathrm{K}(\mathrm{E}^{p^{m}})\) is separable over \(\mathrm{K}(\mathrm{B}_{0})\) , show that for every transcendence basis B of E over K there exists an integer \(n \geqslant 0\) such that \(\mathrm{K}(\mathrm{E}^{p^{n}})\) is separable over \(\mathrm{K}(\mathrm{B})\) .
\item
  Deduce that if the condition in a) holds, E admits a separating transcendence basis over K (if S is a p-basis of E over \(\mathrm{K}(\mathrm{E}^{p})\) and B a transcendence basis of E over K containing S (V, p.~176, Ex. 6), show that B is a separating transcendence basis of E over K, using Ex. 6 of V, p.~176).
\item
  Let K be a field of characteristic p \textgreater{} 0 and x an element transcendental over K (in an algebraically closed extension \(\Omega\) of K). Show that the union E of the extensions \(\mathrm{K}(x^{p^{-n}})\) of K is a separable extension of K such that tr. \(\deg_{K}E = 1\) but which does not admit a separating transcendence basis.
\item
  Let K be a perfect field. Show that the perfect closure of the field of rational fractions \(\mathrm{K}(\mathrm{T})\) is an extension of K of transcendence degree 1 which does not admit a separating transcendence basis.
\end{enumerate}

\begin{enumerate}
\def\labelenumi{\arabic{enumi})}
\setcounter{enumi}{1}
\item
  Let K be field of characteristic p \textgreater{} 0 and E an extension of finite transcendence degree over K. For E to admit a separating transcendence basis over K it is necessary and sufficient that E should be separable over K and the degree of imperfection of E over K equal to its transcendence degree over K.
\item
  Let \(\mathbf{E}\) be an extension of a field \(\mathbf{K}\) and \(\mathbf{F}\) an extension of \(\mathbf{E}\) .
\end{enumerate}

\begin{enumerate}
\def\labelenumi{\alph{enumi})}
\item
  If E admits a separating transcendence basis over K and if F admits a separating transcendence basis over E, then F admits a separating transcendence basis over K.
\item
  If F admits a separating transcendence basis over K and if tr. \(\deg_{K}E<+\infty\) , then E admits a separating transcendence basis over K (reduce to the case where tr. \(\deg_{K}F<+\infty\) and apply Ex. 1).
\item
  If F admits a finite separating transcendence basis over K and is separable over E then F admits a separating transcendence basis over E (use Ex. 1).
\end{enumerate}

\begin{enumerate}
\def\labelenumi{\arabic{enumi})}
\setcounter{enumi}{3}
\item
  Let K be a field of characteristic p \textgreater{} 0 whose absolute degree of imperfection (V, p.~170, Ex. 1) is equal to 1. Show that for an extension E of K to be separable over K it is necessary and sufficient that E should contain no element p-radical over K and not belonging to K (if \(a \in K\) forms a p-basis of K over \(K^{p}\) show that \(\mathbf{K}^{p^{-1}} = \mathbf{K}(a^{p^{-1}})\) and E are linearly disjoint over K).
\item
  \begin{enumerate}
  \def\labelenumii{\alph{enumii})}
  \tightlist
  \item
    Let K be a field of characteristic p \textgreater{} 0 and E an extension of K admitting a separating transcendence basis over K. Show that the intersection L of the fields \(\mathrm{K}(E^{p^{n}})\) , where n runs over N, is an algebraic extension of K. (Consider first the case where \(\operatorname{tr} \cdot \deg_{K} E < +\infty\) ; show that L is relatively perfect over K, hence E is separable over L. Now use Ex. 2 and Ex. 6 of V, p.~176. In the general case, if B is a separating transcendence basis of E over K, for each \(x \in L\) , there is a finite subset \(B_{0}\) of B such that the minimal polynomial of x over each of \(\mathrm{K}(B_{0}^{p^{n}})\) is the same as its minimal polynomial over \(\mathrm{K(B)}\) ; now deduce that if \(\mathrm{F} = \mathrm{K}(B_{0})(x)\) , then \(x \in \mathrm{K}(F^{p^{n}})\) for all n.)
  \end{enumerate}
\end{enumerate}

\begin{enumerate}
\def\labelenumi{\alph{enumi})}
\setcounter{enumi}{1}
\tightlist
\item
  Show that the largest perfect field \(E^{p^{\infty}}\) (intersection of the \(E^{p^{n}}\) ) contained in E is an algebraic extension of \(K^{p^{\infty}}\) (use a) and Ex. 8 of V, p.~176). Give an example where E is a finitely generated extension of a perfect field K but \(E^{p^{\infty}}\) is not contained in K.
\end{enumerate}

\begin{enumerate}
\def\labelenumi{\arabic{enumi})}
\setcounter{enumi}{5}
\tightlist
\item
  Let \(\mathbf{K}\) be a field of characteristic \(p > 0\) , \(\mathbf{E}\) an extension of \(\mathbf{K}\) such that there exists an integer \(h \geqslant 0\) for which \(\mathbf{K}(\mathrm{E}^{p^h})\) is separable over \(\mathbf{K}\) (which is always the case when \(\mathbf{E}\) is a finitely generated extension of \(\mathbf{K}\) ).
\end{enumerate}

\begin{enumerate}
\def\labelenumi{\alph{enumi})}
\item
  If C is a \(p\) -basis of E over \(\mathbf{K}(\mathbf{E}^p)\) , there is a subset B of C such that \(\mathbf{B}^{p^h}\) is a \(p\) -basis of \(\mathbf{K}(\mathbf{E}^{p^h})\) over \(\mathbf{K}(\mathbf{E}^{p^{h+1}})\) ; B is an algebraically free set over K. Show that the extension \(\mathbf{F} = \mathbf{K}(\mathbf{E}^{p^h})(\mathbf{B})\) is separable over K and E is an algebraic extension of F.
\item
  Show that B is a \(p\) -basis of F over \(\mathbf{K}(\mathbf{F}^p)\) . Show that if \(x \in \mathbf{E}\) and if \(m\) is the least integer such that \(x^{p^m} \in \mathbf{F}\) , then \(x^{p^m} \in \mathbf{K}(\mathbf{F}^{p^m})\) ; therefore E is a subfield of the field \(\mathbf{K}^{p^{-\infty}}(\mathbf{F})\) (isomorphic to \(\mathbf{K}^{p^{-\infty}} \otimes_{\mathbf{K}} \mathbf{K}\) ). A subextension F of E is said to be distinguished if it is separable over K and if E is a subfield of \(\mathbf{K}^{p^{-\infty}}(\mathbf{F})\) ; F is then a maximal separable extension of K in E.
\item
  Let P be a perfect field of characteristic p \textgreater{} 0, \(K = P(X, Y)\) the field of rational fractions in two indeterminates and \(\Omega\) an algebraically closed extension of the field of rational fractions \(\mathbf{K}(\mathbf{Z})\) over K. If \(u \in \Omega\) is such that \(u^{p} = X + YZ^{p}\) , show that \(\mathrm{E} = \mathrm{K}(\mathrm{Z}, u)\) is an extension of K such that \(\mathrm{K}(\mathrm{Z})\) and \(\mathrm{K}(u)\) are two distinguished subextensions of E, so that there exists no larger separable subextension of E. If \(v \in \Omega\) is such that \(v^{p^{2}} = X + YZ^{p}\) , show that in the extension \(\mathrm{E}' = \mathrm{K}(\mathrm{Z}, v)\) , \(\mathrm{K}(\mathrm{Z})\) is a maximal separable subextension of E but is not a distinguished subextension.
\end{enumerate}

\begin{enumerate}
\def\labelenumi{\arabic{enumi})}
\setcounter{enumi}{6}
\tightlist
\item
  Assume that the hypotheses of Ex. 6 hold, and further that the degree of imperfection of E over K (V, p.~170, Ex. 1) is finite (which is always the case if E is a finitely generated extension of K).
\end{enumerate}

\begin{enumerate}
\def\labelenumi{\alph{enumi})}
\tightlist
\item
  Let B be a \(p\) -basis of F over \(\mathbf{K}(\mathbf{F}^p)\) which is also a separating transcendence basis of F over K (and a transcendence basis of E over K). Show that for every integer \(k > 0\) , \(\mathbf{B}^{p^k}\) is a \(p\) -independent subset of \(\mathbf{K}(\mathbf{E}^{p^k})\) over \(\mathbf{K}(\mathbf{E}^{p^{k+1}})\) . Let \(q = \operatorname{tr} \cdot \deg_{\mathbf{K}} \mathbf{E}\) ; if \(m_k + q\) is the degree of imperfection of \(\mathbf{K}(\mathbf{E}^{p^k})\) over \(\mathbf{K}(\mathbf{E}^{p^{k+1}})\) , show that \(p^f\) , where \(f = \sum_{k=0}^{h-1} m_k\) , is equal to the degree {[}E : F{]} and is therefore independent of the choice of the
\end{enumerate}

distinguished subextension F of E. The integer \(p^{f}\) is called the order of inseparability of E over K and is written \([E:K]_{i}\) ; this number coincides with the inseparable degree denoted by \([E:K]_{i}\) in V, p.46 when E is of finite degree.

\begin{enumerate}
\def\labelenumi{\alph{enumi})}
\setcounter{enumi}{1}
\item
  For \(0 \leqslant k \leqslant h - 1\) , \([\mathbf{K}^{p^{-k}}(\mathbf{E}) : \mathbf{K}^{p^{-k}}]_i\) is equal to \(p^{f_k}\) , where \(f_k = \sum_{j=k}^{h-1} m_j\) .
\item
  Let \(L\) be a subextension of \(E\) , separable over \(K\) and such that \(E\) is \(p\) -radical and of finite degree over \(L\) ; show that \([E:L] \geqslant [E:K]_i\) (if \([L(E^{p^k}):K(E^{p^k})] = p^{r_k}\) , show that \(r_{k+1} \leqslant r_k + q\) , by considering a basis of \(L(E^{p^k})\) over \(K(E^{p^k})\) and a basis of \(L\) over \(K(L^p)\) ).
\item
  Let \(K'\) be an extension of K such that \(K'\) and E are contained in the same algebraically closed extension \(\Omega\) of K. Show that if \(\mathrm{K}'(\mathrm{F})\) is separable over \(K'\) , it is a distinguished subextension of the extension \(\mathrm{K}'(\mathrm{E})\) of \(K'\) . This is so when \(K'\) and E are algebraically disjoint over K and then we have \([\mathrm{K}'(\mathrm{E}):\mathrm{K}']_{i}\leqslant[\mathrm{E}:\mathrm{K}]_{i}\) ; equality holds when \(K'\) and E are linearly disjoint over K or when \(K'\) is a separable algebraic extension of K. e) Show that \([\mathrm{E}:\mathrm{K}]_{i}\) is the least value of \([\mathrm{E}:\mathrm{K}(\mathrm{B})]_{i}\) for all transcendence bases B of E over K; if \(\bar{K}\) is the relative algebraic closure of K in \(\Omega\) (which is an algebraic closure of K), then \([\mathrm{E}:\mathrm{K}(\mathrm{B})]_{i}=[\mathrm{E}:\mathrm{K}]_{i}[\bar{\mathrm{K}}(\mathrm{E}):\bar{\mathrm{K}}(\mathrm{B})]_{i}\) . If L is a subextension of E such that E is an algebraic extension of L and \(\bar{\mathrm{K}}(\mathrm{E})\) a separable (algebraic) extension of \(\bar{\mathrm{K}}(\mathrm{L})\) , then the largest separable extension F of L contained in E is distinguished.
\item
  Suppose that E is a finitely generated extension of K; show that for every subextension L of E we have \([L:K]_{i} \leqslant [E:K]_{i} \leqslant [L:K]_{i}[E:L]_{i}\) . Give an example where
\end{enumerate}

\[
[ \mathrm{L}: \mathrm{K} ] _ {i} = [ \mathrm{E}: \mathrm{K} ] _ {i} \quad \text { and } \quad [ \mathrm{E}: \mathrm{L} ] _ {i} > 1.
\]

\begin{enumerate}
\def\labelenumi{\arabic{enumi})}
\setcounter{enumi}{7}
\tightlist
\item
  Let \(k\) be a field and \(L\) an extension of \(k\) . Show that the following conditions are equivalent:
\end{enumerate}

\begin{enumerate}
\def\labelenumi{(\roman{enumi})}
\item
  L is a separable algebraic extension of a finitely generated pure transcendental extension.
\item
  \(L\) is a separable extension of \(k\) and
\end{enumerate}

\[
[ \Omega_ {L / k}: L ] = \operatorname{tr} \cdot \deg_ {k} L <   + \infty .
\]

(To prove (ii) \(\Rightarrow\) (i) we may proceed as follows: let \(x_{1},\ldots ,x_{n}\) in L be such that \(dx_{1},\ldots ,dx_{n}\) form a basis of \(\Omega_k(\mathbf{L})\) and put \(\mathbf{K} = k(x_1,\dots,x_n)\) . One first shows that \(\operatorname {tr.deg}_k(\mathbf{K}) = n\) and hence that K is a pure transcendental extension of \(k\) . Let \(\mathbf{L}'\) be a finite algebraic extension of K contained in L. By comparing the canonical mapping \(u:\Omega_k(\mathbf{K})\otimes_\mathbf{K}\mathbf{L}'\to \Omega_k(\mathbf{L}')\) with the canonical mapping \(\Omega_k(\mathbf{K})\otimes_\mathbf{K}\mathbf{L}\to \Omega_k(\mathbf{L})\) , show that \(u\) is injective. Calculate \([\Omega_k(\mathbf{L}'):\mathbf{L}']\) (V, p.~134, Cor. 1). Deduce that \(u\) is bijective and hence that \(\Omega_{\mathbf{K}}(\mathbf{L}') = 0\) . Show that \(\mathbf{L}'\) is separable over K; deduce that L is separable over K.)

\begin{enumerate}
\def\labelenumi{\arabic{enumi})}
\setcounter{enumi}{8}
\tightlist
\item
  With the notation of Ex. 8 let \(\mathbf{M} \subset \mathbf{L}\) be a subextension. Put \(m = \operatorname{tr} \cdot \deg_k \mathbf{M}\) , \(n = \operatorname{tr} \cdot \deg_k \mathbf{L}\) , \(r = n - m\) . Let \(x_1, \ldots, x_n\) in \(\mathbf{L}\) be such that \(\mathbf{L}\) is a separable algebraic extension of \(k(x_1, \ldots, x_n) = \mathbf{K}\) .
\end{enumerate}

\begin{enumerate}
\def\labelenumi{\alph{enumi})}
\item
  Show that, after renumbering the \(x_{i}\) if necessary, we may suppose that \(L\) is algebraic over \(M(x_{1},\ldots ,x_{r})\) .
\item
  Show that there exists a finitely generated subextension \(\mathbf{N} \subset \mathbf{M}\) such that:
\end{enumerate}

\(\alpha)\) L is algebraic over \(\mathbf{N}(x_1,\dots,x_s)\) ;

\(\beta)\) \(\mathbf{K}' = \mathbf{K}(\mathbf{N}(x_1,\dots,x_r))\) is a finite extension of \(\mathbf{N}(x_1,\dots,x_r)\) ;

\(\gamma)\) the canonical homomorphism \(\mathbf{K}^{\prime}\otimes_{\mathbf{N}(x_1,\ldots ,x_r)}\mathbf{M}(x_1,\dots,x_r)\to \mathbf{K}(\mathbf{M}(x_1,\dots,x_r)) = \mathbf{K}''\) is bijective.

\begin{enumerate}
\def\labelenumi{\alph{enumi})}
\setcounter{enumi}{2}
\tightlist
\item
  Let N be a subextension as in b). Show successively that \(K''\) is a separable algebraic extension of \(K'\) (because L is separable over K), that \(M(x_{1}, \ldots, x_{r})\) is a separable algebraic extension of \(N(x_{1}, \ldots, x_{r})\) (use b), \(\gamma)\) and that M is a separable algebraic extension of N. d) Show that M is a separable algebraic extension of a pure transcendental extension (use c) and the fact that N is a finitely generated separable extension of k).
\end{enumerate}

\BourbakiExerciseGroup

\begin{enumerate}
\def\labelenumi{\arabic{enumi})}
\tightlist
\item
  Let \(K_{0}\) be a field of characteristic p \textgreater{} 0, \(\mathbf{K} = \mathbf{K}_{0}(\mathbf{X}, \mathbf{Y})\) the field of rational fractions in two indeterminates over \(K_{0}\) , U and V two indeterminates, \(\Omega\) an algebraically closed extension of \(\mathbf{K}(\mathbf{U}, \mathbf{V})\) , \(\mathbf{E} = \mathbf{K}(\mathbf{U}, u)\) , where \(u^{p} = X + YU^{p}\) and \(F = K(V, v)\) , where \(v^{p} = X + YV^{p}\) ; E and F are not separable over K (V, p.178, Ex. 6, c)).
\end{enumerate}

\begin{enumerate}
\def\labelenumi{\alph{enumi})}
\item
  Show that \(\mathbf{K}\) is relatively algebraically closed in \(\mathbf{E}\) and in \(\mathbf{F}\) (if \(x \in \mathbf{E}\) is algebraic over \(\mathbf{K}\) , show that we must have \(x^p \in \mathbf{K}\) ; if \(x \notin \mathbf{K}\) , then \(\mathbf{E} = \mathbf{K}(\mathbf{U}, x)\) , deduce that then \(\mathbf{X}^{1/p}\) and \(\mathbf{Y}^{1/p}\) would lie in \(\mathbf{K}(x)\) ).
\item
  Show that E and F are linearly disjoint over K but that K is not relatively algebraically closed in \(\mathbf{K}(\mathbf{E} \cup \mathbf{F})\) . (Prove that \(X^{1/p} \in \mathbf{K}(\mathbf{E} \cup \mathbf{F})\) ; deduce that v cannot belong to \(\mathbf{E}(\mathbf{V})\) , and conclude from this that E and F are linearly disjoint over K.)
\end{enumerate}

\begin{enumerate}
\def\labelenumi{\arabic{enumi})}
\setcounter{enumi}{1}
\tightlist
\item
  Let \(\mathbf{K}\) be a field of characteristic \(p > 0\) , \(\Omega\) an extension of \(\mathbf{K}\) and \(\mathbf{E}\) , \(\mathbf{F}\) two subextensions of \(\Omega\) , algebraically disjoint over \(\mathbf{K}\) .
\end{enumerate}

\begin{enumerate}
\def\labelenumi{\alph{enumi})}
\tightlist
\item
  Let \(F_{s}\) be the relative separable closure of K in F. Show that \(E(F_{s})\) is the relative separable closure of E in \(E(F)\) . (Reduce to the case where \(F_{s}=K\) ; let B be a transcendence basis of E over K and \(K'\) the relative algebraic closure of K in F, which is radical over F; then (V, p.~137) \(K'(B)\) is the relative algebraic closure of \(K(B)\) in \(F(B)\) and is p-radical over \(K(B)\) . If \(x\in E(F)\) is algebraic over E, there exists an integer \(r\geqslant0\) such that \(x^{p'}\in M\) , where M is a separable extension of finite degree of \(F(B)\) having a basis \((u_{i})_{1\leqslant i\leqslant n}\) over \(\mathrm{F(B)}\) consisting of elements of E. If \(y = \sum_{i}b_{i}u_{i}\) with \(b_{i}\in \mathrm{F(B)}\) , show
\end{enumerate}

by considering the F(B)-isomorphisms of M into an algebraic closure of E(F) that the \(b_{i}\) belong to K'(B) and deduce that y is p-radical over E.)

\begin{enumerate}
\def\labelenumi{\alph{enumi})}
\setcounter{enumi}{1}
\tightlist
\item
  Suppose that E and F are linearly disjoint over K, that K is relatively algebraically closed in F and that E is a separable extension of K. Show that E is relatively algebraically closed in E(F). (Reduce to the case where E is a separable algebraic extension of K, with the help of V, p.~138. Deduce from a) that the relative algebraic closure of E in E(F) is p-radical over E and then use the fact that if \((a_{\lambda})\) is a basis of E over K, the \(a_{\lambda}^{p'}\) also form a basis of E over K.)
\end{enumerate}

\begin{enumerate}
\def\labelenumi{\arabic{enumi})}
\setcounter{enumi}{2}
\tightlist
\item
  Let E be a field, G a group of automorphisms of E and K the field of invariants of G. a) Show that if there exists no normal subgroup of finite index in G other than G, then E is a regular extension of K.
\end{enumerate}

\begin{enumerate}
\def\labelenumi{\alph{enumi})}
\setcounter{enumi}{1}
\tightlist
\item
  Deduce from a) that if G is isomorphic to a finite product of infinite simple groups, then E is a regular extension of K.
\end{enumerate}

\begin{enumerate}
\def\labelenumi{\arabic{enumi})}
\setcounter{enumi}{3}
\tightlist
\item
  Let K be a field and E, F two extensions of K.
\end{enumerate}

\begin{enumerate}
\def\labelenumi{\alph{enumi})}
\item
  Suppose that K is relatively algebraically closed in F. Show that if \((L, u, v)\) , \((L', u', v')\) are two composite extensions of E and F such that \(u(E)\) and \(v(F)\) (resp \(u'(E)\) and \(v'(F)\) ) are algebraically disjoint over K, then these two composite extensions are isomorphic. (Consider first the case where E is a pure transcendental extension of K; then use Ex. 2, a), as well as Ex. 6, b) of V, p.~154)
\item
  Let \(E_{s}\) , \(F_{s}\) be the relative separable closures of K in E and F respectively. Suppose that all the composite extensions of \(E_{s}\) and \(F_{s}\) are isomorphic. Show that if \((L, u, v)\) and \((L', u', v')\) are two composite extensions of E and F such that \(u(E)\) and \(v(F)\) (resp. \(u'(E)\) and \(v'(F)\) ) are algebraically disjoint over K, then these two composite extensions are isomorphic (use a) and Ex. 11 of V, p.~173).
\end{enumerate}

\begin{enumerate}
\def\labelenumi{\arabic{enumi})}
\setcounter{enumi}{4}
\tightlist
\item
  Let K be a field and p its characteristic exponent. For every integer n \textgreater{} 0 and prime to p, denote by \(\varphi_{\mathrm{K}}(n)\) the degree over K of \(\mathrm{R}_{n}(\mathrm{K}) = \mathrm{K}(\mu_{n})\) , where \(\mu_{n}\) denotes the group of n-th roots of unity in an algebraic closure of K. For example \(\varphi_{\mathrm{Q}}(n) = \varphi(n)\) (V, p.~84, Th. 2); if K is finite with q elements, \(\varphi_{\mathrm{K}}(n)\) is the order of the residue class of q in \((\mathbf{Z}/n\mathbf{Z})^{*}\) (V, p.~97, Cor.); if K is algebraically closed, \(\varphi_{\mathrm{K}}(n) = 1\) for all n.
\end{enumerate}

\begin{enumerate}
\def\labelenumi{\alph{enumi})}
\item
  Let E be the relative algebraic closure of the prime field \(K_{0}\) in K. Show that \(\varphi_{\mathrm{K}}(n) = \varphi_{\mathrm{E}}(n)\) .
\item
  Suppose that K is a finitely generated extension of \(K_{0}\) . Let m be an integer \textgreater0; then there exists an integer N \textgreater0 divisible by every n such that \(\varphi_{\mathrm{K}}(n) \leqslant m\) .
\item
  Let \(s\) be an element of \(\mathrm{GL}_m(\mathbf{K})\) of finite order. Show that \(s^N\) has order a power of \(p\) . (We have \(s = tu = ut\) , where \(u\) has order a power of \(p\) and \(t\) has order \(n\) prime to \(p\) . The element \(t\) is root of a separable polynomial \(\mathbf{X}^n - 1\) as well as of its characteristic polynomial; deduce that the subalgebra \(\mathbf{A}\) of \(\mathbf{M}_m(\mathbf{K})\) generated by \(t\) is etale, of degree \(\leqslant m\) , isomorphic to a product of extensions \(\mathbf{R}_d(\mathbf{K})\) , where \(d\) runs over certain divisors of \(n\) including \(n\) itself. Therefore \(\varphi_{\mathbf{K}}(n) \leqslant m\) ; use \(b\) ).)
\end{enumerate}
